\subsection{Algebras over Separably Closed Fields}

Lemma 1. --- Let D be a field and Z its center. Denote the characteristic exponent of Z by p.~Suppose that for every element a of D, there exists an integer \(m \geqslant 0\) such that \(a^{p^{m}}\) belongs to Z. Then the field D is commutative.

If p = 1, then D = Z. We therefore assume p \textgreater{} 1.

Let us give a proof by contradiction and suppose that D is not commutative. Let \(a\) be an element of \(\mathrm{D} - \mathrm{Z}\) and \(q\) a power of \(p\) such that \(a^q\) belongs to Z. Denote the identity mapping on D by I and the inner automorphism \(x \mapsto axa^{-1}\) from D to D associated with \(a\) by \(\sigma\) ; we have \(\sigma^q = \mathrm{I}\) because \(a^q\) belongs to Z. We have \(\sigma - \mathrm{I} \neq 0\) because \(a\) does not belong to Z, and we have \((\sigma - \mathrm{I})^q = \sigma^q - \mathrm{I} = 0\) because Z has characteristic \(p\) . Let \(f\) be the greatest natural number such that \((\sigma - \mathrm{I})^f \neq 0\) ; we have \(f \geqslant 1\) . Choose an element \(c\) of D such that \((\sigma - \mathrm{I})^f(c) \neq 0\) , and set

\[
x = (\sigma - \mathrm{I}) ^ {f - 1} (c), \qquad y = (\sigma - \mathrm{I}) (x) = (\sigma - \mathrm{I}) ^ {f} (c).
\]

By construction, we have \(y \neq 0\) and \(\sigma(y) = y\) ; if we set \(z = y^{-1}x\) , then it follows that

\[
\sigma (z) = \sigma (y) ^ {- 1} \sigma (x) = y ^ {- 1} (y + x) = 1 + z
\]

and therefore \(\sigma(z^{p^j}) = 1 + z^{p^j}\) for every natural number \(j\) . Choose an integer \(m \geqslant 0\) such that \(z^{p^m}\) belongs to the center Z of D; we have

\[
z ^ {p ^ {m}} = a z ^ {p ^ {m}} a ^ {- 1} = \sigma (z ^ {p ^ {m}}) = 1 + z ^ {p ^ {m}}.
\]

This contradiction implies Lemma 1.

PROPOSITION 3. --- Let K be a separably closed field (V, §7, No.~8, p.~45, Definition 4), and let D be an algebra of finite degree over K that is a field. Then D is commutative.

Denote the characteristic exponent of K by p.~Let a be an element of D. The ring \(K[a]\) is an algebraic extension of K (V, §3, No.~1, p.~17, Corollary 1). Since the field K is separably closed, it follows from V, §7, No.~7, p.~44, Proposition 13 that the algebra \(K[a]\) is a p-radical extension of K. Hence there exists an integer \(m \geqslant 0\) such that \(a^{p^{m}}\) belongs to K. By Lemma 1, the field D is consequently commutative.

COROLLARY. --- Let K be a separably closed field and A a semisimple algebra of finite degree over K. Then there exist an integer \(r \geqslant 0\) , strictly positive integers \(n_{1}, \ldots, n_{r}\) , and extensions \(K_{1}, \ldots, K_{r}\) of K of finite degree such that A is isomorphic to the algebra \(\prod_{i=1}^{r} \mathbf{M}_{n_{i}}(\mathbf{K}_{i})\) .

By the structure theorem for semisimple algebras (VIII, p.~135, Theorem 1), A is isomorphic to an algebra \(\prod_{i=1}^{r} \mathbf{M}_{n_i}(\mathrm{D}_i)\) , where \(r\) is an integer \(\geqslant 0\) , \(n_1, \ldots, n_r\) are strictly positive integers, and \(\mathrm{D}_1, \ldots, \mathrm{D}_r\) are K-algebras of finite degree that are fields. Since the field K is separably closed, each field \(\mathrm{D}_i\) is commutative by Proposition 3; the corollary follows.

